\section{Settling the Claims: The Bit-Sliced Commitment}
\label{sec:commitment}

Both halves of the proof --- the linear PIOP (\Cref{sec:piop}) and the Hadamard
discharge (\Cref{sec:hadamard}) --- terminate in the same place: a batch of
multilinear evaluation claims at a common point $r_0$, asked of the committed
columns under the monomial and Lagrange readings $\psia,\psiz$ (one evaluation
at $\alpha$, two bases). This section summarises the hash-based $\Ftwo[X]$
commitment that certifies them. It is deliberately compact; the full
construction, the four verifier checks, and their soundness proofs are developed
in the companion note \emph{``A Hash-Based Polynomial Commitment Scheme over
$\Ftwo[X]$: the bit-sliced opening''}. What matters here is the one feature the
previous sections were built to exploit: a single bound object per column from
which both readings are read off.

\subsection{Interface}
\label{sec:pcs-interface}

The commitment is a Brakedown/Ligero-style interactive oracle proof of proximity
(IOPP) compiled with a Merkle tree:
\begin{itemize}[leftmargin=2em]
  \item \textsc{Commit}: given the primary witness columns
    $w_1,\dots,w_K\in\cellring^{\,N_{\mathrm{rows}}}$, output one hash root (and
    retain the encoded data as opening hint).
  \item \textsc{Open}: given a point $r_0$, a reading's weight vector, and the
    claimed values, prove that the committed columns' projected MLEs take those
    values at $r_0$.
\end{itemize}
The host protocol calls \textsc{Open} once, at $r_0$, and reads off it both the
$\psia$ claims (linear part) and the $\psiz$ claims (Hadamard discharge).

\subsection{Code and commitment}
\label{sec:pcs-code}

Each column is arranged as a matrix and its rows encoded by a rate-$1/\rho$
$\Ftwo$-linear Repeat--Accumulate--Accumulate code. Two properties carry the
construction (both proved in the companion note): the code \emph{does not widen}
$\cellring$ --- codewords are again bit-polynomials, committable without
growth --- and its $\Ftwo$-linearity \emph{extends scalar-wise} to
$\K$-coefficient rows, so the commit-time codewords double as the proximity
oracle for the opening with no re-encoding. The codewords of all columns at a
given position are hashed into one Merkle leaf, so a single authentication path
reveals every column's value there. The root is the commitment.

\begin{remark}[Provenance, and the proof-size lever]
\label{rem:blaze}
The $\Ftwo$-RAA code and its linear-time, non-widening commit are those of Blaze
\cite{blaze}, a hash-based multilinear commitment over binary fields. What this
paper adds on top is the \emph{bit-sliced} opening (\Cref{sec:bit-slice-fold}):
where Blaze and the IOPPs it composes settle a single evaluation, binding the
width-$\wbit$ fold $a'$ lets the two readings $\psia,\psiz$ share one opening.
Blaze settles its opening not by a flat proximity test but by
\emph{code-switching} to a Reed--Solomon codeword discharged by an inner
foldable IOPP (BaseFold/WHIR) \cite{basefold,whir}, yielding a polylogarithmic
verifier; the deployed system here uses the flat Brakedown test \cite{brakedown},
whose $t$ column queries dominate proof size. That code-switch is orthogonal to
the cell axis --- the $\wbit$ bit positions are carried as passive shared width,
never folded inside the recursion --- so it would shrink the proof while
preserving the dual read-off. This is the recorded direction for closing the
proof-size gap.
\end{remark}

\subsection{The bit-slice fold, and the dual read-off}
\label{sec:bit-slice-fold}

The opening's central object is, per column, the \emph{bit-slice fold}: the
degree-$<\wbit$ univariate over $\K$
\[
  a'_g = \sum_{b=0}^{\wbit-1} a'_{g,b}\,X^b \in \K[X]_{<\wbit},
  \qquad
  a'_{g,b} = \sum_{x\in\{0,1\}^{\nu}} \eq_x(r_0)\,(w_{g})_b(x)\in\K,
\]
whose $b$-th coefficient is the $b$-th bit slice of column $g$, folded along
$\eq(\cdot;r_0)$. Because every reading $\psi_\xi$ (\Cref{def:projection}) is the
inner product of the coefficient vector with $\xi$, \emph{any} projected MLE
evaluation at $r_0$ is a single length-$\wbit$ read-off of $a'_g$:
\[
  \psi_\xi(a'_g) = \sum_b a'_{g,b}\,\xi_b = \mle{\psi_\xi(w_g)}(r_0).
\]
In particular $\psia(a'_g) = \mle{\psia(w_g)}(r_0)$ recovers the linear claim and
$\psiz(a'_g) = \mle{\psiz(w_g)}(r_0)$ the discharge claim, from the \emph{one}
object $a'_g$. This is exactly the dual read-off of \Cref{sec:had-binding}: the
opening binds the full bit-slice fold once, and the two readings are two
length-$\wbit$ inner products against it.

\begin{remark}[Why ``bit-sliced'', and what it buys]
\label{rem:why-bitsliced}
The opening keeps the $\eq$-weights in $\K$ and folds the $\{0,1\}$ bit slices,
rather than lifting the weights into $\Ftwo[X]$ and folding cells. The bound
claim is then a clean degree-$<\wbit$ polynomial over $\K$ --- fixed size $\wbit$
field elements, independent of $\log|\K|$ --- and retains \emph{all} bit slices
rather than one collapsed value. That is what makes one opening serve both
$\psia$ and $\psiz$ (and any further $\Ftwo$-linear functional of the cells): the
nonlinear constraints ride the linear part's opening at the cost of one extra
reading.
\end{remark}

\subsection{The opening, in brief}
\label{sec:pcs-checks}

To certify the $\gamma$-batched claim $a' = \sum_g\gamma_g a'_g$, the opening
runs four checks (full statements and proofs in the companion note):
\begin{enumerate}[leftmargin=2em]
  \item \emph{Evaluation consistency} --- $a'$ is the $\eq(\cdot;r_0)$ fold of the
    prover's row-fold vector;
  \item \emph{Projection discharge} --- $\psi_\xi(a')$ equals the
    $\gamma$-batched claimed values (at $\xi = \alpha$ for the linear part and,
    via \eqref{eq:psi-z-binding}, at the $\psiz$ weights for the discharge);
  \item \emph{Coherence} --- the proximity-tested combined row is consistent with
    the row-fold vector; and
  \item \emph{Encoding consistency} --- at $t$ transcript-chosen codeword
    positions, the committed cells (revealed by Merkle paths) encode the combined
    row.
\end{enumerate}
Check (4) is the proximity test and dominates the cost and the soundness; checks
(1)--(3) are $\K[X]_{<\wbit}$-arithmetic identities binding $a'$ to the committed
data. The opening's knowledge-soundness error is
\[
  \epsilon_{\mathrm{open}} \le
  \epsilon_{\mathrm{prox}}(t,\delta) + O\!\left(\tfrac{1}{|\K|}\right),
\]
with $\epsilon_{\mathrm{prox}}(t,\delta)$ the code's $t$-query proximity error.

\subsection{What it certifies}
\label{sec:pcs-certifies}

An accepting opening extracts committed columns $\widehat w_1,\dots,\widehat w_K$
and certifies that the claimed $\psia$- and $\psiz$-evaluations at $r_0$ are the
true projected evaluations of those columns. Composed back through
\Cref{sec:piop,sec:hadamard}, this certifies that the extracted witness satisfies
the entire UAIR --- every linear constraint (via the $\psia$ claims) and every
Hadamard relation (via the $\psiz$ claims and the binding). The next section
assembles these into the soundness of the whole argument.
